\section{Related Work}
\label{sec:related}

% todo:
% http://www.eecs.qmul.ac.uk/~gretay/papers/onward2017.pdf

\paragraph{Superoptimization for LLVM and GCC}
%
A superoptimizer by Sands~\cite{Sands11} pointed out many opportunities
for optimizations that had not yet been implemented in LLVM around 2010
and 2011.
%
This tool had many similarities to \tool{}: it extracted directed
acyclic graphs of LLVM IR during compilation and then attempted to
optimize them, it focused on the integer subset of LLVM, and it
presented its results ordered by static profile count.
%
On the other hand, Sands' tool relied on testing to perform
equivalence checking and was, therefore, unsound; this did not lead to
miscompilations since it worked offline, and was only intended to
generate suggestions for developers to implement.
%
To synthesize right-hand sides, this superoptimizer applied various
heuristic simplifications of the left-hand side.
%
Analogously, the GNU superoptimizer~\cite{Granlund92} was an unsound,
enumeration-based tool that was used to improve GCC's code generation
by deriving ways to turn code with control flow into straight-line
code.


\paragraph{An online superoptimizer}
%
Most previous superoptimizers have been intended for offline use: it
has not been practical to run them during regular compiles.
%
Bansal and Aiken's tool~\cite{Bansal06} is one of the few exceptions:
it achieved good performance by caching its results in a database
(among other techniques).
%
Additionally, it was sound, using a SAT solver to verify equivalence.
%
This tool was a direct inspiration for \tool{}, though our work
improves upon it in important ways, for example by using synthesis
instead of enumeration to construct RHSs, by learning dataflow facts
from diverging and converging control flow, and by extracting
instructions that are related by dataflow rather than instructions
that happen to be close to each other in the instruction stream.
%
Bansal and Aiken's tool, on the other hand, could deal with memory
accesses and vector instructions, while \tool{} cannot.


\paragraph{The original superoptimizer}
%
Massalin~\cite{Massalin87} coined the term \emph{superoptimizer} to
emphasize the fact that regular optimizers produce code that is
far from optimal.
%
His tool was offline and unsound, using enumeration to discover
right-hand sides and testing to rule out inequivalent ones.


\paragraph{The earliest superoptimizers}
%
Fraser's 1979 paper~\cite{Fraser79} is the earliest work we know of
that captures all of the essential ideas of superoptimization:
extracting instruction sequences from real programs, searching for
cheaper ways to compute them, and checking each candidate using an
equivalence checker.
%
Moreover, unlike much subsequent work, Fraser's equivalence checker
was sound.
%
Subsequently, in 1984, Kessler~\cite{Kessler84} developed a sound
superoptimizer for a Lisp compiler.


\paragraph{Recent superoptimizers}
%
There has been significant progress in superoptimizers in recent
years.
%
Optgen~\cite{Buchwald15}, like \tool{} and unlike almost all of the
other work described in this section, operates on IR rather than
assembly language.
%
It generates all optimizations up to a cost limit, rather than
extracting code sequences from programs.
%
Unlike \tool{}, it can synthesize optimizations containing symbolic
constants.
%
STOKE~\cite{Schkufza13,Schkufza14,Sharma15} is a sound superoptimizer
that uses randomized search, rather than solver-based synthesis, to
find cheap right-hand sides.
%
It can handle many features that \tool{} cannot, including memory,
floating point instructions, approximations, and even, in conjunction
with the DDEC tool~\cite{Sharma13}, loops.
%
Chlorophyll~\cite{Phothilimthana14} is a synthesis-aided compiler for
a 144-core chip.
%
Finally, GreenThumb~\cite{Phothilimthana16} is a generic
superoptimizer that can be easily retargeted; it has been used to
optimize LLVM IR\@.




\paragraph{Synthesis}
%
There has been enormous progress in program sketching and synthesis in
the last few years.
%
Our work mainly builds upon a single paper, Gulwani et
al.~\cite{Gulwani11}.
